% Options for packages loaded elsewhere
\PassOptionsToPackage{unicode}{hyperref}
\PassOptionsToPackage{hyphens}{url}
\PassOptionsToPackage{space}{xeCJK}
%
\documentclass[
]{ctexart}
\usepackage{amsmath,amssymb}
\usepackage{iftex}
\ifPDFTeX
  \usepackage[T1]{fontenc}
  \usepackage[utf8]{inputenc}
  \usepackage{textcomp} % provide euro and other symbols
\else % if luatex or xetex
  \usepackage{unicode-math} % this also loads fontspec
  \defaultfontfeatures{Scale=MatchLowercase}
  \defaultfontfeatures[\rmfamily]{Ligatures=TeX,Scale=1}
\fi
\usepackage{lmodern}
\ifPDFTeX\else
  % xetex/luatex font selection
  \ifXeTeX
    \usepackage{xeCJK}
    \setCJKmainfont[]{Songti SC}
          \fi
  \ifLuaTeX
    \usepackage[]{luatexja-fontspec}
    \setmainjfont[]{Songti SC}
  \fi
\fi
% Use upquote if available, for straight quotes in verbatim environments
\IfFileExists{upquote.sty}{\usepackage{upquote}}{}
\IfFileExists{microtype.sty}{% use microtype if available
  \usepackage[]{microtype}
  \UseMicrotypeSet[protrusion]{basicmath} % disable protrusion for tt fonts
}{}
\makeatletter
\@ifundefined{KOMAClassName}{% if non-KOMA class
  \IfFileExists{parskip.sty}{%
    \usepackage{parskip}
  }{% else
    \setlength{\parindent}{0pt}
    \setlength{\parskip}{6pt plus 2pt minus 1pt}}
}{% if KOMA class
  \KOMAoptions{parskip=half}}
\makeatother
\usepackage{xcolor}
\usepackage[margin=2.2cm]{geometry}
\setlength{\emergencystretch}{3em} % prevent overfull lines
\providecommand{\tightlist}{%
  \setlength{\itemsep}{0pt}\setlength{\parskip}{0pt}}
\setcounter{secnumdepth}{5}
\renewcommand{\and}{\\}
\ifLuaTeX
  \usepackage{selnolig}  % disable illegal ligatures
\fi
\IfFileExists{bookmark.sty}{\usepackage{bookmark}}{\usepackage{hyperref}}
\IfFileExists{xurl.sty}{\usepackage{xurl}}{} % add URL line breaks if available
\urlstyle{same}
\hypersetup{
  pdftitle={Survival Analysis},
  pdfauthor={Rui},
  hidelinks,
  pdfcreator={LaTeX via pandoc}}

\title{Survival Analysis}
\author{Rui}
\date{\texttt{r\ Sys.Date()}}

\begin{document}
\maketitle

\renewcommand*\contentsname{目录}
{
\setcounter{tocdepth}{3}
\tableofcontents
}
\hypertarget{packages}{%
\section{Packages}\label{packages}}

\begin{itemize}
\tightlist
\item
  Survival Analysis

  \begin{itemize}
  \tightlist
  \item
    survival
  \item
    survminer
  \end{itemize}
\item
  Plots of Survival Analysis

  \begin{itemize}
  \tightlist
  \item
    ggfortify
  \item
    ggplot
  \end{itemize}
\end{itemize}

\texttt{\{r\ message=FALSE,results=\textquotesingle{}hide\textquotesingle{},warning=FALSE\}\ library(survival)\ library(survminer)\ library(ggplot2)\ library(ggfortify)}

Dataset ``lung'' in package ``survival'' is of interest, which contains
the information about survival in patients with advanced lung cancer
from the North Central Cancer Treatment Group. Performance scores rate
how well the patient can perform usual daily activities. Variables
include:

\begin{itemize}
\tightlist
\item
  inst: Institution code
\item
  time: Survival time in days
\item
  status: censoring status \texttt{1=censored}, \texttt{2=dead}
\item
  age: Age in years
\item
  sex: \texttt{Male=1} \texttt{Female=2}
\item
  ph.ecog: ECOG performance score as rated by the physician.
  \texttt{0=asymptomatic}, \texttt{1=\ symptomatic} but completely
  ambulatory, \texttt{2=\ in\ bed\ \textless{}50\%\ of\ the\ day},
  \texttt{3=\ in\ bed\ \textgreater{}\ 50\%\ of\ the\ day\ but\ not\ bedbound},
  \texttt{4\ =\ bedbound}
\item
  ph.karno: Karnofsky performance score
  (\texttt{bad=0}-\texttt{good=100}) rated by physician
\item
  pat.karno: Karnofsky performance score as rated by patient
\item
  meal.cal: Calories consumed at meals
\item
  wt.loss: Weight loss in last six months (pounds)
\end{itemize}

\texttt{\{r\ warning=FALSE,message=FALSE\}\ attach(lung)} \# Survival
Object

\texttt{survival::Surv()} creates a survival object, with censored data
marked with ``\texttt{+}'' following up censoring time and death cases
simply showing the death time. The result is usually used as a response
variable in a model formula.

Two most important parameters:

\begin{itemize}
\tightlist
\item
  \texttt{time} :For right censored data, this is the follow up time.
  For interval data, the first argument is the starting time for the
  interval.
\item
  \texttt{event}: The status indicator, normally \texttt{0=alive},
  \texttt{1=dead}. Other choices are \texttt{TRUE/FALSE}
  (\texttt{TRUE\ =\ death}) or \texttt{1/2} (\texttt{2=death}). For
  interval censored data, the status indicator is 0=right censored,
  1=event at time, 2=left censored, 3=interval censored. For multiple
  endpoint data the event variable will be a factor, whose first level
  is treated as censoring. Although unusual, the event indicator can be
  omitted, in which case all subjects are assumed to have an event.
\end{itemize}

\texttt{\{r\}\ survobj\ =\ Surv(time,\ status==2)}

\hypertarget{surv-model}{%
\section{Surv Model}\label{surv-model}}

For the overall survival analysis, just specify the null hypothesis with
\texttt{Surv\textasciitilde{}1}, i.e., NO groups.

\texttt{\{r\}\ sfit\ \textless{}-\ survfit(Surv(time,\ status==2)\textasciitilde{}1,\ data=lung)}

The fitted survival model itself will tell you the results, with
observations, events, median and 95\% CI.

\texttt{\{r\}\ sfit}

If more details are needed, i.e., the instantaneous events, and
estimates of survival, std.err, 95\% CI, just \texttt{summary(sfit)}.

Moreover, if \texttt{time} is designated, then the \texttt{summary()}
function will return the estimated results.

\texttt{\{r\}\ summary(sfit,time\ =\ 200)}

However, if group is of interest, specify it in the formula with
\texttt{Surv\textasciitilde{}group}.

\texttt{\{r\}\ sfit\_sex\ =\ survfit(Surv(time,status==2)\textasciitilde{}strata(sex))\ sfit\_sex}

From the model itself, we can get an overall picture of the differences
between groups.

\texttt{\{r\}\ summary(sfit\_sex,time\ =\ 365)}

It's ideal that the 95\% CI of groups doesn't intersect, and a
significant result can be obtained instantly. However, that's not the
case, take a step further and use Cox regression to see exactly.

\hypertarget{cox-regression}{%
\section{Cox Regression}\label{cox-regression}}

\texttt{\{r\}\ coxph(Surv(time,\ status==2)\textasciitilde{}sex)} where
``\texttt{exp(coef)}'' is the HR (hazard ratio)， which measures the
relative risk between two groups.

Theoretically, it's better only open to binary variables in Cox
regression. However, for categorical variables with more than 2 levels,
use ``\texttt{as.factor(cat\_var)}'' with ``\texttt{model\ =\ TRUE}'' to
generate corresponding dummy variables. For continuous variables,
they'll be included into the regression model as usual, but the model
will be plotted with expected value of continuou variables.

However, similar to \texttt{coxph()} which returns the differences
through the view of regression, the simpler version of such comparison
can be accomplished by ``\texttt{survdiff()}'', also in package
``survival''. Both of the functions have the same formula.

From ``\texttt{survdiff()}'', we can get the intermediate products when
computing the log rank, or the chi-square test. However, this wouldn't
return the coefficient.

Jointly speaking, the combined results of two functions emphasize the
fact that the log rank test (named Mantel-Haenszel test), is a
chi-square test in essence.

\texttt{\{r\}\ surv\_diff\ =\ survdiff(Surv(time,\ status==2)\textasciitilde{}sex)\ surv\_diff\ surv\_diff\$chisq}

\hypertarget{surv-plots}{%
\section{Surv Plots}\label{surv-plots}}

Finally, we can visualize the survival analysis.

Note that even if analysed data has been attached, if the
``\texttt{data}'' isn't passed into \texttt{survfit()}, then an error
will be raised.

The parameter ``\texttt{risk.table}'' provides you with the freedom to
decide whether or not to add a risk table below the Kaplan-Meier
curve(s). Meanwhile, various options are available!

\texttt{\{r\ warning=FALSE\}\ ggsurvplot(sfit\_sex,data\ =\ lung,\ \ \ \ \ \ \ \ \ \ \ \ conf.int=TRUE,\ pval=TRUE,\ \ \ \ \ \ \ \ \ \ \ \ \ risk.table=\textquotesingle{}abs\_pct\textquotesingle{},\ \ \ \ \ \ \ \ \ \ \ \ legend.labs=c("Male",\ "Female"),\ legend.title="Sex",\ \ \ \ \ \ \ \ \ \ \ \ title="Kaplan-Meier\ Curve\ for\ Lung\ Cancer\ Survival")}

Moreover, the parameter ``\texttt{fun}'' is quite useful, which will
define a transformation of the survival curve.

\begin{itemize}
\tightlist
\item
  ``pct'': for survival probability in percentage, \textbf{by default}.
\item
  ``event'': plots cumulative events, \(f(y)=1-y\).
\item
  ``cumhaz'': plots the cumulative hazard function, \(f(y)=-\log(y)\).
\end{itemize}

\texttt{\{r\}\ ggsurvplot(sfit\_sex,data\ =\ lung,\ \ \ \ \ \ \ \ \ \ \ \ conf.int=TRUE,\ pval=TRUE,\ \ \ \ \ \ \ \ \ \ \ \ \ \#\ risk.table=\textquotesingle{}abs\_pct\textquotesingle{},\ \ \ \ \ \ \ \ \ \ \ \ legend.labs=c("Male",\ "Female"),\ legend.title="Sex",\ \ \ \ \ \ \ \ \ \ \ \ \#\ title="Kaplan-Meier\ Curve\ for\ Lung\ Cancer\ Survival",\ \ \ \ \ \ \ \ \ \ \ \ fun\ =\ \textquotesingle{}event\textquotesingle{})}

\texttt{\{r\}\ ggsurvplot(sfit\_sex,data\ =\ lung,\ \ \ \ \ \ \ \ \ \ \ \ conf.int=TRUE,\ pval=TRUE,\ \ \ \ \ \ \ \ \ \ \ \ \ \#\ risk.table=\textquotesingle{}abs\_pct\textquotesingle{},\ \ \ \ \ \ \ \ \ \ \ \ legend.labs=c("Male",\ "Female"),\ legend.title="Sex",\ \ \ \ \ \ \ \ \ \ \ \ title="Kaplan-Meier\ Curve\ for\ Lung\ Cancer\ Survival",\ \ \ \ \ \ \ \ \ \ \ \ fun\ =\ \textquotesingle{}cumhaz\textquotesingle{})}

\end{document}
